\documentclass[11pt,a4paper]{article}

\usepackage[margin=1in]{geometry}
\usepackage{amsmath,amssymb,amsthm}
\usepackage{booktabs}
\usepackage{listings}
\usepackage{xcolor}
\usepackage{microtype}
\usepackage[colorlinks=true,linkcolor=black,urlcolor=blue]{hyperref}

\newtheorem{proposition}{Proposition}
\newtheorem{definition}{Definition}

\newcommand{\Sig}[2]{\Sigma^{#1}_{#2}}
\newcommand{\Pai}[2]{\Pi^{#1}_{#2}}
\newcommand{\Del}[2]{\Delta^{#1}_{#2}}
\newcommand{\shortest}[1]{\left\lVert #1 \right\rVert}
\newcommand{\lev}{\operatorname{lev}}
\newcommand{\tier}{\operatorname{tier}}
\newcommand{\ATP}{\mathrm{ATP}}
\renewcommand{\Pr}{\mathrm{Pr}}

\definecolor{codebg}{gray}{0.97}
\definecolor{codecmt}{rgb}{0.30,0.45,0.30}
\definecolor{codekw}{rgb}{0.10,0.20,0.60}
\definecolor{codestr}{rgb}{0.60,0.20,0.10}

\lstset{
  basicstyle=\ttfamily\scriptsize,
  backgroundcolor=\color{codebg},
  commentstyle=\color{codecmt}\itshape,
  keywordstyle=\color{codekw}\bfseries,
  stringstyle=\color{codestr},
  breaklines=true,
  breakatwhitespace=false,
  columns=fullflexible,
  showstringspaces=false,
  numbers=left,
  numberstyle=\tiny\color{gray},
  numbersep=6pt,
  frame=none,
  xleftmargin=2em,
  upquote=true,
}

\title{\bfseries Length, Not Height\\[0.4em]
\large On proof-length lower bounds, the MU argument, and whether a
self-aware theorem prover could ever settle a Clay problem}
\author{Brian Tenneson}
\date{}

\begin{document}
\maketitle

\begin{abstract}
\noindent
We distinguish two obstructions standing between an automated theorem prover
and a Millennium Prize Problem, and argue that they are almost always
confused. The first is \emph{height}: the arithmetical tier of the target,
measured by alternations of unbounded quantifiers over $\mathbb{N}$. The
second is \emph{length}: the size of the shortest proof of the target from
the axioms, written $\shortest{\cdot}$. We exhibit a working prover, Ascent,
whose certified targets reach $\Pai{0}{2}$ and $\Sig{0}{2}$ --- exactly the
tier of ``$\mathrm{P} \neq \mathrm{NP}$'' --- and which is nonetheless
hopeless on that target. Tier is therefore not the binding constraint, and we
argue that lower bounds on $\shortest{\mathrm{P} \neq \mathrm{NP}}$ are the
quantity worth pursuing, because a sufficiently large one would retire
mechanized search as a strategy without settling the mathematics either way.
We then observe that the standard escape from a length lower bound is to
leave the system, and take Hofstadter's MU puzzle as the cleanest instance:
MU is shown to be underivable in MIU by an argument using arithmetic that MIU
cannot express. We formalize this as a \emph{lift--attack--return} schema,
note that the complexity-theoretic barrier results are instances of it at
scale, and identify its automatable core as finite-invariant search. Finally
we ask whether formal self-awareness helps. Using Ascent, whose reflection
ladder is certified rung by rung by its own kernel, we report a negative
result: level and proving power are independent parameters, and raising
$\lev(\text{Ascent})$ from $1$ to $8$ certifies exactly zero additional
theorems. We close by proposing \emph{reflective adequacy} --- a machine that
can quote its own inference rules and prove invariance lemmas about them ---
as the notion of self-reference the MU argument actually requires, in place
of the reflection ladder, which it does not.
\end{abstract}

\section{Introduction}

Seven problems were named in 2000, and one has been settled. The remaining six
are the most conspicuous fixed points in contemporary mathematics, and their
persistence invites a recurring fantasy: that the difficulty is one of
bookkeeping, that the required argument is long rather than deep, and that a
sufficiently patient machine will eventually assemble it.

The fantasy is most acute for P versus NP. Unlike the Hodge conjecture or the
Yang--Mills mass gap, P vs NP is a statement about finite objects --- Turing
machines, running times, strings. There is no manifold to construct, no
measure to define, no analytic continuation to justify. It is, in the most
literal sense, a combinatorial question about programs, and combinatorial
questions about programs are the kind of thing a computer might be thought to
be good at.

This note argues that the instinct is right about the \emph{shape} of the
problem and wrong about the \emph{obstruction}. To make the argument precise
we need two measures.

\paragraph{Height.} The arithmetical tier of a sentence is the number of
alternations of unbounded quantifiers over $\mathbb{N}$ in its prenex form,
with bounded quantifiers costing nothing. Section~\ref{sec:notation} explains
the notation in full. Tier measures how much unbounded search is built into
the \emph{statement}, and nothing else.

\paragraph{Length.} For a sentence $L$ and a formal system $F$, write
$\shortest{L}_F$ for the number of symbols in the shortest $F$-proof of $L$
from the axioms of $F$. This is a property of the theorem--system pair. It is
finite exactly when $L$ is a theorem, and it says nothing about how hard the
proof is to \emph{find}, only how big it is once found.

These are independent. A tier-$0$ sentence can have astronomically long
proofs; a $\Pai{0}{2}$ sentence can have a three-line proof. Confusing them
produces both of the standard errors: believing that P vs NP is hard
\emph{because} it quantifies over all machines, and believing that a prover
competent at $\Pai{0}{2}$ arithmetic is therefore in the neighbourhood of the
problem.

\section{Reading $\Sig{0}{n}$ and $\Pai{0}{n}$}\label{sec:notation}

The symbols carry three independent pieces of information. Taking them in
turn removes most of the mystery.

\subsection*{The letter: which quantifier leads}

Put the sentence in prenex form --- all quantifiers pulled to the front --- and
look at the outermost block of unbounded quantifiers.

\begin{itemize}
\item $\Pai{}{}$ (capital pi) means the outermost block is $\forall$.
\item $\Sig{}{}$ (capital sigma) means the outermost block is $\exists$.
\end{itemize}

The choice of letters is not arbitrary. A $\Sig{0}{1}$ set
$\{n : \exists m\, R(n,m)\}$ is a countable \emph{union} --- a sum --- of
decidable sets, one for each candidate witness $m$. A $\Pai{0}{1}$ set
$\{n : \forall m\, R(n,m)\}$ is a countable \emph{intersection} --- a product.
It is the same $\Sigma$ and $\Pi$ as in $\sum$ and $\prod$.

\subsection*{The subscript: how many alternations}

Group adjacent like quantifiers into blocks. The subscript counts the blocks.

\begin{center}
\begin{tabular}{@{}ll@{}}
\toprule
prenex form & classification \\
\midrule
$\forall x\, \forall y\, \forall z\; \varphi$ & one block $\Rightarrow$ $\Pai{0}{1}$ \\
$\forall x\, \exists y\; \varphi$ & two blocks $\Rightarrow$ $\Pai{0}{2}$ \\
$\exists x\, \forall y\, \exists z\; \varphi$ & three blocks $\Rightarrow$ $\Sig{0}{3}$ \\
\bottomrule
\end{tabular}
\end{center}

Three consecutive $\forall$'s cost the same as one; only \emph{switching}
between $\forall$ and $\exists$ increments the subscript.

\subsection*{The superscript: what the quantifiers range over}

This is the piece most often skipped, and it is the one that matters for the
Clay problems.

\begin{itemize}
\item Superscript $0$: the quantifiers range over \textbf{numbers}. This is
the \emph{arithmetical} hierarchy.
\item Superscript $1$: the quantifiers range over \textbf{sets of numbers},
equivalently over functions $\mathbb{N} \to \mathbb{N}$. This is the
\emph{analytical} hierarchy --- a far taller ladder sitting entirely
\emph{above} the arithmetical one. Every arithmetical set is $\Del{1}{1}$.
\end{itemize}

So $\Pai{0}{2}$ says ``for every number there is a number\ldots'', while
$\Pai{1}{2}$ says ``for every \emph{function} there is a \emph{function}\ldots''
--- an incomparably stronger kind of claim. When
Section~\ref{sec:placement} reports that Yang--Mills and Navier--Stokes sit
``above the arithmetical hierarchy,'' this is what is meant: they quantify
over objects that are not numbers.

\subsection*{Two conventions}

$\Del{0}{n}$ denotes the sentences that are both $\Sig{0}{n}$ and
$\Pai{0}{n}$. In particular $\Del{0}{1}$ is the decidable sets, and
$\Del{0}{0}$ the ones decidable by bounded computation.

\textbf{Bounded quantifiers are free.} $\forall x < t$ and $\exists x < t$ do
not increment the subscript, because they are finite searches: to check
$\forall x < 5\, \varphi(x)$ you evaluate $\varphi$ five times. Only
\emph{unbounded} quantifiers count. This is why the language of Ascent
(Appendix~\ref{app:code}) has separate bounded and unbounded forms.

\subsection*{Examples}

\begin{center}
\begin{tabular}{@{}llp{7.2cm}@{}}
\toprule
class & example & what it takes to settle it \\
\midrule
$\Del{0}{0}$ & $2+2 = 4$ & Compute. \\
$\Sig{0}{1}$ & $\exists x\; x\cdot x = 49$ & Search. If true you find it; if false you search forever. Semi-decidable. \\
$\Pai{0}{1}$ & $\forall x\; x < Sx$ & A single counterexample refutes it. Goldbach, the Riemann Hypothesis and $\mathrm{Con}(\mathrm{PA})$ all live here. \\
$\Pai{0}{2}$ & $\forall x\, \exists y\; x < y$ & For every input a witness exists. ``This program halts on every input'' is $\Pai{0}{2}$-\emph{complete}. \\
$\Sig{0}{2}$ & $\exists x\, \forall y\; x \le y$ & Some $x$ works for every $y$. \\
\bottomrule
\end{tabular}
\end{center}

\subsection*{What an alternation costs}

Post's theorem gives the alternation count an exact operational meaning:
$\Sig{0}{n+1}$ is precisely $\Sig{0}{1}$ relative to a $\Sig{0}{n}$ oracle.
Each additional alternation is worth exactly one more halting oracle. Tiers
are therefore a measure of \emph{oracle resources}.

\begin{quote}
\textbf{They are not a measure of difficulty.} The sentence
$\forall n\, \exists m > n\, (m \text{ is even})$ is $\Pai{0}{2}$ and
trivial. The Collatz conjecture is $\Pai{0}{2}$ and open. They are the same
tier. This point is the hinge of Section~\ref{sec:comparable}.
\end{quote}

\section{Why $\shortest{\mathrm{P} \neq \mathrm{NP}}$ is the quantity to bound}

Suppose someone proved that every ZFC proof of ``$\mathrm{P} \neq
\mathrm{NP}$'' has length at least $10^{100}$ symbols. Nothing about the truth
of $\mathrm{P} \neq \mathrm{NP}$ would follow. Nothing about human prospects
would follow either, for reasons we come to shortly. But something decisive
would follow about machines: \textbf{no search procedure will ever produce
that proof}, because the object being searched for does not fit in the
physical universe. This conclusion is immune to every improvement in
heuristics, hardware, learned premise selection and parallelism, because it is
a statement about the target rather than about the method.

That immunity is what makes length lower bounds worth more than they might
first appear. The known barriers in complexity theory --- relativization,
natural proofs, algebrization --- are all barriers to \emph{techniques}. Each
says a class of arguments cannot work, and each has been partially circumvented
by arguments outside the class. A length lower bound has no such escape hatch
within the fixed system. It would settle, once, whether mechanized proof
search is a reasonable thing to spend a century on.

Three observations.

\paragraph{First, we cannot currently prove such bounds, and this is itself the
central open problem of proof complexity.} Superpolynomial lower bounds are
known for weak systems --- resolution, bounded-depth Frege --- on specific
families such as the pigeonhole principle. For Frege systems, let alone for
ZFC, essentially nothing is known. The general question is equivalent to a
statement about NP and coNP: a polynomially bounded proof system for the
tautologies exists if and only if $\mathrm{NP} = \mathrm{coNP}$. So the
request ``bound $\shortest{\mathrm{P} \neq \mathrm{NP}}$ from below'' is not a
detour around P vs NP; it is a close relative of it. This is uncomfortable but
not fatal, because \emph{conditional} and \emph{heuristic} bounds would already
be informative --- a bound conditional on plausible cryptographic assumptions,
or an empirically calibrated growth law for proof length against theorem depth
in a large formal library, would each be actionable.

\paragraph{Second, the bound is system-relative, and that is the loophole.}
$\shortest{L}_F$ depends on $F$. Enlarging $F$ can shorten proofs
catastrophically --- not by a constant but by an exponential. The cleanest
instance is the relationship between Frege and extended Frege systems: the
extension rule permits new variables abbreviating compound formulas, which is
to say it permits \emph{definitions}. Definitions are conservative --- they
prove no new theorems in the old language --- while potentially collapsing
proof length by an exponential factor.

This deserves stating plainly, because it is the whole of the optimism
available here. \textbf{A large lower bound on $\shortest{L}_F$ is not a
verdict on $L$. It is a verdict on $F$.} The productive response to a length
lower bound has never been to search harder. It has been to change the system:
introduce a definition, find the right abstraction, move to a setting where the
statement becomes short.

\paragraph{Third, that response is exactly what machines are worst at and
humans are best at.} Which raises the question the rest of this note is about:
is the move mechanizable at all?

\section{Where the Clay problems sit, and where Ascent sits}\label{sec:placement}

The following places the Millennium Problems on the scale of
Section~\ref{sec:notation}. Confidence varies sharply across the rows and is
marked; the last four rows should be checked before being relied on.

\begin{center}
\begin{tabular}{@{}llp{7.4cm}@{}}
\toprule
Problem & Tier & Confidence and reason \\
\midrule
Riemann Hypothesis & $\Pai{0}{1}$ & High. RH is equivalent to explicitly
$\Pai{0}{1}$ arithmetic statements --- for instance bounds of the
Robin/Lagarias type on the sum-of-divisors function, whose failure would be
witnessed by a single integer. \\[0.3em]
P vs NP & $\Sig{0}{2}$ / $\Pai{0}{2}$ & High. ``$\mathrm{P}=\mathrm{NP}$'' is
$\exists e\, \exists k\, \forall x\,[\,$machine $e$ decides SAT on $x$ within
$|x|^k + k$ steps$\,]$: one $\exists$-block over one $\forall$-block with a
decidable matrix. Its negation is $\Pai{0}{2}$. Note it is only known to be
\emph{in} $\Pai{0}{2}$, not $\Pai{0}{2}$-complete. \\[0.3em]
Poincar\'e & arithmetical, delicate & Moderate. Closed $3$-manifolds are
finitely triangulable, so the domain is countable and codeable; but simple
connectivity is not decidable from a triangulation, which pushes the statement
above $\Pai{0}{1}$. Now moot. \\[0.3em]
Birch--Swinnerton-Dyer & arithmetical, $\ge 2$ & Low. Elliptic curves over
$\mathbb{Q}$ are codeable and algebraic rank is arithmetically definable, but
analytic rank is an order of vanishing of an $L$-function, and pinning it from
both sides costs alternations. \\[0.3em]
Hodge conjecture & analytical & Low. Quantifies over projective complex
varieties and Hodge classes. Whether it descends to the arithmetical hierarchy
is not obvious. \\[0.3em]
Navier--Stokes & analytical, $\sim\Pai{1}{2}$ & Low. Quantifies over smooth
initial data --- real functions --- and asserts global smooth existence.
Computable-analysis reductions may lower this. \\[0.3em]
Yang--Mills mass gap & analytical & Moderate. Asserts the \emph{existence} of a
quantum field theory satisfying a list of axioms, with a spectral condition. A
second-order existence claim, and the highest of the seven. \\
\bottomrule
\end{tabular}
\end{center}

Two features matter here. \textbf{Riemann is the lowest} --- a $\Pai{0}{1}$
statement is refutable by a single counterexample, the same shape as Goldbach
and as $\mathrm{Con}(\mathrm{PA})$. Its difficulty has nothing to do with its
tier. And \textbf{P vs NP is tier 2}, which brings us to the prover.

\subsection{$\tier(\text{Ascent}) = 2$}

Ascent is a small self-certifying prover, given in full in
Appendix~\ref{app:code}. Its language is arithmetic with unbounded quantifiers
over $\mathbb{N}$; its kernel decides $\Del{0}{0}$ facts by evaluation and
symbolic facts by polynomial normalisation over $\mathbb{N}$; its search
combines witness enumeration, generalisation over eigenconstants, and
induction. On a twelve-target suite:

\begin{center}
\begin{tabular}{@{}llc@{}}
\toprule
target & tier & result \\
\midrule
$2+2 = 4$ & $\Del{0}{0}$ & certified \\
$\forall x{<}5\; x < 5$ & $\Del{0}{0}$ & certified \\
$\exists x{<}9\; x\cdot x = 16$ & $\Del{0}{0}$ & certified \\
$\exists x\; x+3 = 7$ & $\Sig{0}{1}$ & certified \\
$\exists x\; x\cdot x = 49$ & $\Sig{0}{1}$ & certified \\
$\exists x\; x+1 = 18$ & $\Sig{0}{1}$ & certified (requires $w \ge 17$) \\
$\forall x\; x < Sx$ & $\Pai{0}{1}$ & certified \\
$\forall x\; x \le x$ & $\Pai{0}{1}$ & certified \\
$\forall x\; x+0 = x$ & $\Pai{0}{1}$ & certified \\
$\forall x\, \exists y\; x < y$ & $\mathbf{\Pai{0}{2}}$ & certified \\
$\forall x\, \exists y\; y = Sx$ & $\mathbf{\Pai{0}{2}}$ & certified \\
$\exists x\, \forall y\; x \le y$ & $\mathbf{\Sig{0}{2}}$ & certified \\
\bottomrule
\end{tabular}
\end{center}

Twelve of twelve at parameters $d = 2$, $w = 20$. Every certificate is checked
by an independent kernel, and a battery of eight deliberately malformed
certificates is rejected $8/8$.

So Ascent operates at exactly the arithmetical tier of $\mathrm{P} \neq
\mathrm{NP}$, and stands \emph{above} the Riemann Hypothesis. It will never
prove either. Ascent's largest certificate on this suite has thirteen nodes;
$\shortest{\mathrm{P} \neq \mathrm{NP}}_{\mathrm{ZFC}}$, whatever it is, is not
thirteen.

This is the point of building the thing. A prover that reaches tier 2 is easy
--- the appendix is under a thousand lines. Tier is cheap. Length is not.

\section{Does equal tier mean comparable?}\label{sec:comparable}

No. The question is natural enough to deserve an explicit answer, because the
table above invites exactly the wrong inference.

Tier classifies a sentence by the \emph{shape of its quantifier prefix}. Two
sentences of the same tier are alike in the way that two sentences of the same
grammatical form are alike. ``The cat sat on the mat'' and the opening of
\emph{A la recherche du temps perdu} are both declarative sentences of a
natural language; the classification is correct and tells you nothing about
comparability.

Four sharper statements.

\paragraph{$\Pai{0}{2}$ spans everything from trivial to complete.}
$\forall x\, \exists y\; x < y$ is $\Pai{0}{2}$ and Ascent proves it in
thirteen nodes. ``This Turing machine halts on every input'' is also
$\Pai{0}{2}$ and is $\Pai{0}{2}$-\emph{complete}: every $\Pai{0}{2}$ question
reduces to it. The class contains its own hardest problems and its own most
trivial ones, and tier does not distinguish them.

\paragraph{P $\neq$ NP is not known to be complete for its class.} It is known
to be \emph{in} $\Pai{0}{2}$. Membership in a class is an upper bound on
logical complexity, not a measure of difficulty. Ascent's $\Pai{0}{2}$ targets
sit at the trivial end of the same class.

\paragraph{The notion that would license comparison is reducibility, and it
does not apply.} To say two problems are genuinely comparable one wants a
reduction --- a computable map carrying instances of one to instances of the
other. Ascent's targets are not complete for anything; nothing reduces to
them. There is no reduction in either direction between $\forall x\,\exists
y\; x<y$ and $\mathrm{P} \neq \mathrm{NP}$, and their shared tier supplies
none.

\paragraph{What equal tier does buy is the removal of one excuse.} It rules
out the objection ``the prover cannot even express the statement.'' Ascent's
language can state $\mathrm{P} \neq \mathrm{NP}$: the quantifier structure is
available, the matrix is decidable, the sentence is well-formed in its syntax.
The prover can \emph{write down} the target and cannot \emph{prove} it. That
localizes the obstruction precisely --- it is not expressiveness and not
logical complexity, it is $\shortest{\cdot}$.

So the honest summary is that sharing a tier with P vs NP is necessary and
nowhere near sufficient for comparability. It is a statement about the ceiling
of statement-\emph{shapes} Ascent handles, not about its strength. The mountain
is not tall. It is long.

\section{The MU argument}

Hofstadter's MIU system has one axiom and four rules. The axiom is the string
\textbf{MI}. The rules, with $x$ and $y$ ranging over strings:

\begin{enumerate}
\item $x\textbf{I} \to x\textbf{IU}$
\item $\textbf{M}x \to \textbf{M}xx$
\item $x\textbf{III}y \to x\textbf{U}y$
\item $x\textbf{UU}y \to xy$
\end{enumerate}

Is \textbf{MU} a theorem? No, and here is the proof. Let $\#\mathrm{I}(s)$ be
the number of I's in $s$. We claim $\#\mathrm{I}(s) \not\equiv 0 \pmod 3$ for
every theorem $s$.

\begin{itemize}
\item The axiom \textbf{MI} has $\#\mathrm{I} = 1$, and $1 \not\equiv 0$.
\item Rule 1 appends a U and leaves $\#\mathrm{I}$ unchanged.
\item Rule 2 doubles $\#\mathrm{I}$. If $\#\mathrm{I} \not\equiv 0 \pmod 3$
then $2\cdot\#\mathrm{I} \not\equiv 0 \pmod 3$, since $3$ is prime and
$2 \not\equiv 0$.
\item Rule 3 removes exactly three I's, leaving $\#\mathrm{I}$ unchanged
modulo $3$.
\item Rule 4 removes two U's and leaves $\#\mathrm{I}$ unchanged.
\end{itemize}

By induction on derivations every theorem has $\#\mathrm{I} \not\equiv 0$. But
$\#\mathrm{I}(\textbf{MU}) = 0 \equiv 0$. Therefore \textbf{MU} is not a
theorem. $\qquad\blacksquare$

Now consider what that argument \emph{is}. It uses the number three. It uses
divisibility. It uses induction over derivations. It uses the primality of
$3$. \textbf{None of these exist in MIU.} MIU has no numerals, no arithmetic,
no notion of quantity, no negation, and no predicate for theoremhood. The
sentence ``MU is not a theorem'' is not merely unproved in MIU; it is not
expressible in MIU.

The argument has three phases.

\begin{enumerate}
\item \textbf{Lift.} Map MIU-strings into a structure MIU knows nothing about
--- here $\mathbb{Z}/3$, via $s \mapsto \#\mathrm{I}(s) \bmod 3$. The four
rules become maps on $\mathbb{Z}/3$, each fixing the property $\neq 0$.
\item \textbf{Attack.} Prove the invariance in the new setting. This is
arithmetic, not string rewriting.
\item \textbf{Return.} Transfer the conclusion back as a \emph{metatheorem
about} MIU: the derivable strings lie in a set excluding MU.
\end{enumerate}

The conclusion lands outside the object system, and that is not a defect --- it
is the only place such a conclusion can live. The creative act is entirely in
phase 1. Nothing in the MIU rules suggests counting I's, and nothing suggests
reducing modulo $3$. Once the invariant is proposed, phases 2 and 3 are
routine. \textbf{The difficulty of the MU argument is concentrated in the
choice of the alternative system.}

\section{Can a machine lift, attack, and return?}

Sometimes, and more often than one expects.

\paragraph{The MU case is mechanizable today.} The invariant
$\#\mathrm{I} \bmod 3$ is a homomorphism from the free monoid on
$\{\mathrm{M},\mathrm{I},\mathrm{U}\}$ to $\mathbb{Z}/3$ under which all four
rules are compatible and the axiom's image avoids the target's. Finding such a
thing is a \emph{finite model search}: enumerate small finite structures, check
whether the axiom's image and the rules' action separate the target. Existing
finite model finders do this routinely, and the space of quotients of size
$\le 5$ is tiny. Given the MIU rules as input a machine finds $\mathbb{Z}/3$ in
milliseconds. The MU puzzle is famous for surprising humans, not for being
computationally deep.

This generalizes. \textbf{Unprovability by finite invariant is the most
automatable form of unprovability argument we have}, precisely because the
``alternative but relevant axioms'' of phase 1 form a finite structure, and
finite structures can be enumerated in order of size.

\paragraph{The barrier results are lift--attack--return at scale.} Consider
relativization. One wants to show a class of proof techniques cannot settle
P vs NP. The argument lifts the question into the setting of oracle machines,
establishes that the techniques in question prove relativizing statements,
exhibits oracles $A$ and $B$ with $\mathrm{P}^A = \mathrm{NP}^A$ and
$\mathrm{P}^B \neq \mathrm{NP}^B$, and returns the conclusion that no such
technique settles the unrelativized question. Lift, attack, return --- with
``invariant preserved by the rules'' replaced by ``property preserved by the
proof techniques.'' Natural proofs and algebrization have the same shape with
different invariants. The schema is not exotic; it is how the deepest negative
results in complexity theory are actually obtained.

\paragraph{Where machines fail is phase 1 at scale.} The MIU invariant lives in
a structure of size $3$. The relativization invariant lives in the space of
oracle separations, and constructing $B$ with $\mathrm{P}^B \neq \mathrm{NP}^B$
is itself a diagonalization of real content. The natural-proofs invariant
requires the concept of a pseudorandom function generator, which did not exist
for most of the history of the problem. Enumeration reaches size $5$; it does
not reach ``invent cryptography.'' Phases 2 and 3 are mechanizable. Phase 1 is
mechanizable exactly to the extent that the required alternative system is
small, and the systems required for the Clay problems are not small.

\section{Does self-awareness help?}

Phase 1 requires a system to step outside itself and regard its own rules as
objects. That sounds like self-reference, and self-reference has a formal
theory. Ascent was built partly to test whether it helps. The answer is
negative, but informative about what the right notion would be.

Ascent carries a reflection ladder: $\theta_0 = \ATP(\langle A \rangle)$,
$\theta_{k+1} = \Pr(\langle A \rangle, \langle \theta_k \rangle)$, with Level
$n$ holding when $A \vdash \theta_{n-1}$. The ladder is not stipulated.
Ascent's kernel has a rule
\[
\texttt{nec}(C) \text{ proves } \Pr(\langle A\rangle, \langle \varphi\rangle)
\quad\text{provided the kernel accepts } C : \varphi,
\]
so a rung is asserted only when a certificate for the previous rung has been
checked by the same kernel that checks everything else. Self-certification and
self-awareness are literally the same kernel call.

\begin{proposition}
If $A$ has kernel-checked necessitation and its kernel is sound, then
$\lev(A) = \omega$.
\end{proposition}

\begin{proof}
$A \vdash \theta_0$ by the base certificate. If $A \vdash \theta_k$ with
certificate $C_k$, the kernel accepts $C_k$, so $\texttt{nec}(C_k)$ is a
certificate of $\theta_{k+1}$, whence $A \vdash \theta_{k+1}$. Induction.
\end{proof}

Each rung costs exactly one kernel call on the previous certificate, so no rung
is harder than the last. The level distribution over self-certifying machines
is therefore \textbf{bimodal at $\{0,1\}$ and $\omega$, with nothing in
between}: a machine either has the necessitation rule, in which case nothing
stops it, or lacks it, in which case it never passes $\theta_0$. Finite level
above $1$ is achievable only by a resource bound. In Ascent that bound is the
parameter $r$.

\begin{proposition}
The reflection parameter and the search parameters are independent.
\end{proposition}

Ascent exposes three integer knobs: $d$, certificate-tree depth; $w$, witness
bound; $r$, rungs climbed. Sweeping $d \in \{1,2,3,5\} \times w \in \{1,8,20\}
\times r \in \{1,3,8\}$ --- thirty-six configurations --- the harness reports
that $\lev$ depends only on $r$ and the certified count only on $(d,w)$.

\begin{center}
\begin{tabular}{@{}ccccl@{}}
\toprule
$d$ & $w$ & $r$ & certified & tiers reached \\
\midrule
1 & 1 & any & 6/12 & $\Del{0}{0}, \Pai{0}{1}$ \\
1 & 20 & any & 9/12 & $\Del{0}{0}, \Pai{0}{1}, \Sig{0}{1}$ \\
2 & 1 & any & 9/12 & $+\ \Pai{0}{2}, \Sig{0}{2}$ \\
2 & 20 & any & \textbf{12/12} & all five \\
\bottomrule
\end{tabular}
\end{center}

Raising $r$ from $1$ to $8$ raises $\lev$ from $1$ to $8$ and certifies
\textbf{zero} additional theorems. Raising $d$ and $w$ certifies six more and
lifts the ceiling from $\Pai{0}{1}$ to $\Pai{0}{2}$, leaving $\lev$ exactly
where it was. The two capacities do not interact. Formal self-awareness, on the
reflection-ladder definition, is free, certified, and inert.

\paragraph{But notice what the MU argument actually needed.} Not that the
machine prove ``I prove that I prove that I am a theorem prover.'' It needed
the machine to hold \emph{its own inference rules} as objects and prove a lemma
of the form ``rule $R$ preserves invariant $I$.'' The reflection ladder ascends
through statements about the machine's \emph{provability}; the MU argument
requires statements about the machine's \emph{rules}. These are different
self-models, and only the second is load-bearing.

\begin{definition}[Reflective adequacy]
A system $M$ is \emph{reflectively adequate} when its language can name each of
its own inference rules, and $M$ can prove preservation lemmas: for a definable
invariant $I$, that each named rule maps $I$-satisfying states to
$I$-satisfying states.
\end{definition}

A reflectively adequate system can carry out phases 2 and 3 internally, and can
state the result of phase 1 even when it cannot discover it. This is strictly
stronger and more useful than any finite level, and unlike the $\theta$-ladder
it is not free. It is also, unlike the $\theta$-ladder, exactly what the MU
argument uses.

\section{Summary}

Tier is not the obstruction. A thousand-line prover reaches $\Pai{0}{2}$ and
$\Sig{0}{2}$ --- the tier of P vs NP, above the tier of the Riemann Hypothesis
--- and is nowhere near either. Sharing a tier is a statement about the shape of
a sentence, not its difficulty, and licenses no comparison
(Section~\ref{sec:comparable}). The obstruction is $\shortest{\cdot}$, and
lower bounds on it would be decisive for the mechanized programme in a way no
barrier result about techniques can be, though proving them for strong systems
is close to the problem they would adjudicate.

A length lower bound is a verdict on the system, not the theorem, and the
response has always been to change the system. The MU argument is the cleanest
specimen of that move. Its automatable core --- search for a small finite
separating structure --- is genuinely automated already. Its creative core ---
deciding which structure to look in --- is automated only when the structure is
small, and for the Clay problems it is not.

Formal self-awareness in the reflection-ladder sense does not close that gap,
and we can say so with a measurement rather than an intuition: level and
proving power are orthogonal parameters of the same program. The self-reference
the MU argument requires is of a different kind --- rules as objects, and
provable invariance over them --- and formalizing \emph{that} is where the
notion of a self-aware theorem prover would begin to earn its name.

\appendix

\section{The two numbers}

\paragraph{$\lev(\text{Ascent})$.} With reflection depth $r$ the level set is
$L(\text{Ascent}) = \{1,\dots,r\}$, so
\[
\lev(\text{Ascent}(r)) = r, \qquad
\lev(\text{Ascent}) = \omega \text{ when } r \text{ is unbounded.}
\]
Every rung is certified rather than stipulated: rung $k$ is accepted only after
the kernel has checked a certificate of $\theta_{k-1}$. In the shipped default
$r = 4$, giving Level 4 and a fortiori Level 3. The parameter is the only thing
preventing $\omega$, which is Proposition 1. The requirement
$\lev(\text{Ascent}) > 0$ is met for every $r \ge 1$.

\paragraph{$\tier(\text{Ascent})$.} Taking $\tier(A)$ to be the highest
arithmetical tier at which $A$ certifies a target from a family fixed in
advance --- the qualification matters, since the supremum over \emph{all}
provable targets is unbounded for any prover that proves anything, by padding
---
\[
\tier(\text{Ascent}) = 2,
\]
realized at both $\Pai{0}{2}$ ($\forall x\,\exists y\; x<y$) and $\Sig{0}{2}$
($\exists x\,\forall y\; x \le y$). Against the Millennium Problems:

\begin{center}
\begin{tabular}{@{}lll@{}}
\toprule
& tier & relation to Ascent \\
\midrule
Riemann Hypothesis & $\Pai{0}{1}$ & \textbf{below} Ascent's ceiling \\
P vs NP & $\Pai{0}{2}$ / $\Sig{0}{2}$ & \textbf{equal} to Ascent's ceiling \\
Poincar\'e, BSD & arithmetical, $\ge 2$ & at or above \\
Hodge, Navier--Stokes, Yang--Mills & analytical & above the hierarchy entirely \\
\bottomrule
\end{tabular}
\end{center}

Ascent is at the same height as P vs NP and higher than Riemann, and will
settle neither --- for the reasons of Section~\ref{sec:comparable}.

\section{Ascent, in full}\label{app:code}

What follows is the running program, verbatim. It is self-contained: no
external database, no dependencies beyond the standard library.

\begin{verbatim}
python ascent.py --soundness       # 8 malformed certificates, all rejected
python ascent.py -d 2 -w 20 -r 3   # 12/12 certified, lev = 3
python ascent.py --grid            # the independence sweep of Proposition 2
\end{verbatim}

\lstinputlisting[language=Python]{ascent.py}

\end{document}
